% Options for packages loaded elsewhere
\PassOptionsToPackage{unicode}{hyperref}
\PassOptionsToPackage{hyphens}{url}
\documentclass[
  ignorenonframetext,
]{beamer}
\newif\ifbibliography
\usepackage{pgfpages}
\setbeamertemplate{caption}[numbered]
\setbeamertemplate{caption label separator}{: }
\setbeamercolor{caption name}{fg=normal text.fg}
\beamertemplatenavigationsymbolsempty
% remove section numbering
\setbeamertemplate{part page}{
  \centering
  \begin{beamercolorbox}[sep=16pt,center]{part title}
    \usebeamerfont{part title}\insertpart\par
  \end{beamercolorbox}
}
\setbeamertemplate{section page}{
  \centering
  \begin{beamercolorbox}[sep=12pt,center]{section title}
    \usebeamerfont{section title}\insertsection\par
  \end{beamercolorbox}
}
\setbeamertemplate{subsection page}{
  \centering
  \begin{beamercolorbox}[sep=8pt,center]{subsection title}
    \usebeamerfont{subsection title}\insertsubsection\par
  \end{beamercolorbox}
}
% Prevent slide breaks in the middle of a paragraph
\widowpenalties 1 10000
\raggedbottom
\AtBeginPart{
  \frame{\partpage}
}
\AtBeginSection{
  \ifbibliography
  \else
    \frame{\sectionpage}
  \fi
}
\AtBeginSubsection{
  \frame{\subsectionpage}
}
\usepackage{iftex}
\ifPDFTeX
  \usepackage[T1]{fontenc}
  \usepackage[utf8]{inputenc}
  \usepackage{textcomp} % provide euro and other symbols
\else % if luatex or xetex
  \usepackage{unicode-math} % this also loads fontspec
  \defaultfontfeatures{Scale=MatchLowercase}
  \defaultfontfeatures[\rmfamily]{Ligatures=TeX,Scale=1}
\fi
\usepackage{lmodern}
\ifPDFTeX\else
  % xetex/luatex font selection
\fi
% Use upquote if available, for straight quotes in verbatim environments
\IfFileExists{upquote.sty}{\usepackage{upquote}}{}
\IfFileExists{microtype.sty}{% use microtype if available
  \usepackage[]{microtype}
  \UseMicrotypeSet[protrusion]{basicmath} % disable protrusion for tt fonts
}{}
\makeatletter
\@ifundefined{KOMAClassName}{% if non-KOMA class
  \IfFileExists{parskip.sty}{%
    \usepackage{parskip}
  }{% else
    \setlength{\parindent}{0pt}
    \setlength{\parskip}{6pt plus 2pt minus 1pt}}
}{% if KOMA class
  \KOMAoptions{parskip=half}}
\makeatother
\usepackage{longtable,booktabs,array}
\usepackage{caption}
\captionsetup[table]{skip=6pt}
\newcounter{none} % for unnumbered tables
\usepackage{calc} % for calculating minipage widths
\usepackage{caption}
% Make caption package work with longtable
\makeatletter
\def\fnum@table{\tablename~\thetable}
\makeatother
\setlength{\emergencystretch}{3em} % prevent overfull lines
\providecommand{\tightlist}{%
  \setlength{\itemsep}{0pt}\setlength{\parskip}{0pt}}
% Auto-generated by infrastructure/rendering/_pdf_latex_helpers.py::extract_math_font_preamble.
% Minimal header passed to Pandoc via -H so Beamer slides pick up the same math font
% as the combined-PDF path. Adjust the manuscript's preamble.md to change the font globally.
\usepackage{unicode-math}
\setmathfont{latinmodern-math.otf}

% Manuscript macro/package fallbacks (newcommand -> providecommand).
% Core mathematics
\usepackage{amsmath}
\usepackage{amssymb}
% Document layout
% Tables
\usepackage{booktabs}
% Algorithm typesetting (for pseudocode in §2 Methodology)
% Code listings
\usepackage{listings}
% Typography and formatting
% Cross-references and citations
% ── Unicode-capable mono font for code listings ──────────────────────
% JuliaMono (TeX Live 2026) covers the full math/Greek glyph set used in
% scientific code blocks; the default lmmono lacks \alpha/\mu/\partial/\nabla.
%
% The hardcoded TeX Live 2026 path below is portable to most macOS / Linux
% TeX Live installs that placed JuliaMono under the default texmf-dist path.
% If your TeX Live version differs (e.g., 2025, 2027) or your distribution
% installs fonts elsewhere, comment out the explicit Path = line — fontspec
% will then search the system font path. If JuliaMono is not installed at
% all, comment out the whole block; xelatex falls back to lmmono with
% reduced Unicode coverage but the document still compiles.
% Math font for unicode-math: Latin Modern Math (TeX Live) has full BMP
% coverage including U+2223 (\mid), U+226A/226B (\ll/\gg), and the Greek/
% blackboard letters used in equations. Without an explicit \setmathfont,
% unicode-math falls back to lmroman text font which lacks several glyphs
% and emits "Missing character" warnings on every \mid in math mode.

% Slide layout defaults for warning-clean scientific decks.
\usepackage{etoolbox}
\IfFileExists{xurl.sty}{\usepackage{xurl}}{}
\IfFileExists{seqsplit.sty}{\usepackage{seqsplit}}{\newcommand{\seqsplit}[1]{#1}}
\protected\def\breakseq#1{\seqsplit{#1}}
\protected\def\breaktt#1{\begingroup\ttfamily\seqsplit{#1}\endgroup}
\setlength{\emergencystretch}{6em}
\tolerance=5000
\hbadness=10000
\hfuzz=1pt
\setlength{\tabcolsep}{2pt}
\AtBeginEnvironment{longtable}{\tiny\renewcommand{\arraystretch}{0.86}\setlength{\tabcolsep}{1pt}}
\AtBeginEnvironment{tabular}{\tiny\renewcommand{\arraystretch}{0.86}\setlength{\tabcolsep}{1pt}}
\AtBeginEnvironment{equation}{\tiny}
\AtBeginEnvironment{equation*}{\tiny}
\AtBeginEnvironment{align}{\tiny}
\AtBeginEnvironment{align*}{\tiny}
\AtBeginEnvironment{itemize}{\footnotesize}
\AtBeginEnvironment{enumerate}{\footnotesize}
\AtBeginEnvironment{description}{\footnotesize}
\setbeamerfont{caption}{size=\tiny}
\setbeamerfont{caption name}{size=\tiny}
\setbeamerfont{normal text}{size=\small}
\setbeamerfont{frametitle}{size=\small}
\setbeamerfont{section title}{size=\footnotesize}
\setbeamerfont{subsection title}{size=\footnotesize}
\setbeamertemplate{section page}{%
  \centering
  \begin{beamercolorbox}[sep=12pt,center,wd=\paperwidth]{section title}
    \parbox{0.86\paperwidth}{\centering\usebeamerfont{section title}\insertsection\par}
  \end{beamercolorbox}
}
\setbeamertemplate{subsection page}{%
  \centering
  \begin{beamercolorbox}[sep=8pt,center,wd=\paperwidth]{subsection title}
    \parbox{0.86\paperwidth}{\centering\usebeamerfont{subsection title}\insertsubsection\par}
  \end{beamercolorbox}
}
\setlength{\abovecaptionskip}{2pt}
\setlength{\belowcaptionskip}{0pt}

% Natbib and cross-reference fallbacks — slides don't load natbib
% or cleveref, but combined-PDF manuscript prose may emit these
% commands. The fallback renders citations as a bracketed key list
% and cross-references as detokenized labels so slides don't fail on
% undefined control sequences or raw underscores. \providecommand is
% a no-op if packages load later.
\providecommand{\citep}[1]{[#1]}
\providecommand{\citet}[1]{#1}
\providecommand{\citealp}[1]{#1}
\providecommand{\citeauthor}[1]{#1}
\providecommand{\citeyear}[1]{#1}
\providecommand{\cref}[1]{\texttt{\detokenize{#1}}}
\providecommand{\Cref}[1]{\texttt{\detokenize{#1}}}
\providecommand{\crefrange}[2]{\texttt{\detokenize{#1}}--\texttt{\detokenize{#2}}}
\providecommand{\Crefrange}[2]{\texttt{\detokenize{#1}}--\texttt{\detokenize{#2}}}
\providecommand{\paragraph}[1]{\textbf{#1}\ }

\newtheorem{proposition}{Proposition}
\newtheorem{hypothesis}{Hypothesis}
\newtheorem{remark}[theorem]{Remark}
\newtheorem{axiom}{Axiom}
\newtheorem{property}{Property}
\usepackage{bookmark}
\IfFileExists{xurl.sty}{\usepackage{xurl}}{} % add URL line breaks if available
\urlstyle{same}
% fallback for those not using the hyperref driver hyperxmp:
\makeatletter
\@ifundefined{xmpquote}{\newcommand{\xmpquote}[1]{#1}}{}
\makeatother
\hypersetup{
  hidelinks,
  pdfcreator={LaTeX via pandoc}}

\author{\texorpdfstring{}{}}
\date{}

\begin{document}

\section{Boundary Comparators: What Non-Linux and Specialized Systems
Teach}\label{sec:comparators}

\begin{frame}[fragile,allowframebreaks]{Boundary Comparators: What
Non-Linux and Specialized Systems Teach}
The candidates examined so far live inside the Linux workstation and
server lanes. This section places seven systems that sit outside those
lanes --- privacy distributions, a non-Linux BSD, formally verified and
capability-based microkernel systems, a mobile platform, and
offensive-toolkit distributions --- into the same property vocabulary.
Their value here is comparative: each one demonstrates, in a deployed
artifact, a property that the composition argument of this review needs,
and each also demonstrates the boundary at which its claim stops.
Reading them as competitors for a single ``most secure system'' title
misses what they actually teach. All entries follow the evaluation
discipline of sec.~4: documented designs,
explicitly scoped claims, no numeric scores.

\end{frame}

\begin{frame}[fragile,allowframebreaks]{Boundary Comparators: What
Non-Linux and Specialized Systems Teach}
This lane also moved during the review window. OpenBSD shipped two
releases and, more tellingly, used its errata process to \emph{remove} a
previously shipped mitigation; seL4's verification frontier extended to
a new architecture and a new proof class; Sculpt 26.04 reworked how
platform drivers obtain authority; and Tails rebuilt on a new Debian
base while preserving its persistence model across automatic upgrades.
For systems whose value is precisely their engineering discipline,
release mechanics are evidence, not trivia (O. Project 2026b; seL4
Foundation 2026c; Labs 2026; T. Project 2026c).

\end{frame}

\begin{frame}[fragile,allowframebreaks]{Boundary Comparators: What
Non-Linux and Specialized Systems Teach}
\begin{longtable}[]{@{}
  >{\raggedright\arraybackslash}p{(\linewidth - 6\tabcolsep) * \real{0.2500}}
  >{\raggedright\arraybackslash}p{(\linewidth - 6\tabcolsep) * \real{0.2500}}
  >{\raggedright\arraybackslash}p{(\linewidth - 6\tabcolsep) * \real{0.2500}}
  >{\raggedright\arraybackslash}p{(\linewidth - 6\tabcolsep) * \real{0.2500}}@{}}
\caption{Non-Linux and boundary comparators: documented contribution,
boundary or caveat, and the lesson each contributes to the composition
argument. Assessments are analytical judgments from documented designs
and project documentation.}\label{tbl:comparators}\tabularnewline
\toprule\noalign{}
\begin{minipage}[b]{\linewidth}\raggedright
System
\end{minipage} & \begin{minipage}[b]{\linewidth}\raggedright
Documented contribution
\end{minipage} & \begin{minipage}[b]{\linewidth}\raggedright
Boundary or caveat
\end{minipage} & \begin{minipage}[b]{\linewidth}\raggedright
Lesson for the composition argument
\end{minipage} \\
\midrule\noalign{}
\endfirsthead
\toprule\noalign{}
\begin{minipage}[b]{\linewidth}\raggedright
System
\end{minipage} & \begin{minipage}[b]{\linewidth}\raggedright
Documented contribution
\end{minipage} & \begin{minipage}[b]{\linewidth}\raggedright
Boundary or caveat
\end{minipage} & \begin{minipage}[b]{\linewidth}\raggedright
Lesson for the composition argument
\end{minipage} \\
\midrule\noalign{}
\endhead
\textbf{Whonix} & A gateway/workstation split that routes all traffic
through a dedicated Tor gateway; its workstation documentation
explicitly warns that compromise exposes the workstation's credentials
and browser data (W. Project 2026b, 2026c). Whonix 17 has reached end of
security support (with final deprecation notices through early and mid
2026), making Whonix 18 the only supported line (W. Project 2026a). &
Anonymity is not credential protection: the project's own documentation
treats workstation compromise as exposing workstation-held secrets.
Running a deprecated release line to avoid migration is exactly the
maintenance failure the supported-line discipline warns against. &
Boundary placement of networking --- keeping the network-facing
component in a separate trust domain --- is a structural pattern that
composes with compartmentalization; it does not make an authorized agent
harmless. \\
\textbf{Tails} & An amnesic live environment with Tor networking, plus
optional encrypted Persistent Storage whose trade-offs the project
documents (T. Project 2026a, 2026b). The 7.x series is built on the
Debian 13 base, preserves Persistent Storage across its automatic
upgrade path, and current releases (7.12) carry a documented notice
about Secure Boot certificate expiry for users who enabled Secure Boot
on earlier releases (T. Project 2026c). & Amnesia addresses trace
reduction and session reset, not information exposed while the session
is live; enabling persistence deliberately reintroduces state that must
then be protected. Even an amnesic system's boot chain needs certificate
maintenance. & Session-state hygiene is one contribution to
\breaktt{persistence\_recovery}, and it must not be confused with a
recovery plan for credentials used during the session. \\
\textbf{OpenBSD} & Documented secure-by-default service choices,
privilege separation, and exploit mitigations including W\^{}X and
system-wide hardening (O. Project 2026d, 2026a). Releases 7.8
(2025-10-22) and 7.9 (2026-05-19) document the current state: a new
\breaktt{\_\_pledge\_open(2)} interface refined the pledge-to-open path,
and \texttt{pinsyscall} gained a PT\_LOAD refusal check that prevents
mapping crafted executable segments over pinned entry points; kernel
guard pages back the discipline (Project 2025, 2026c). & 2026 errata for
7.8 and 7.9 \emph{removed} the previously documented pledge
\texttt{tmppath} permission and fixed a related namei race --- the
project retracted one of its own mitigations when its maintenance burden
and race behavior outweighed its value (O. Project 2026b). The
documented record also does not establish superiority for a modern
browser-heavy, GPU-dependent, AI-development workstation. & Disciplined
trusted-computing-base design is a property, not a platform identity;
its techniques transfer as ideas even where the OS itself does not fit
the workload. Even the most disciplined base revises its own mitigations
through a public errata process. \\
\textbf{seL4} & Formal verification of a small microkernel with
explicitly scoped proofs and assumptions, including boot, hardware, and
DMA-related qualifications (seL4 Foundation 2026a; Klein et al. 2009).
The 2026 seL4 milestone set documents the seL4 16.0.0 release,
completion of the functional verification of the MCS kernel on 64-bit
RISC-V, and the first confidentiality proof on AArch64, with Microkit
2.3.0 adding x86\_64 IOMMU support for the smaller deployment path (seL4
Foundation 2026c, 2026b). & The proof boundary does not extend to the
browser, driver stack, firmware chain, or user workflow that a real
deployment adds on top. New proofs on new architectures broaden the
frontier; they do not lift the existing scope caveats. & Strong
assurance of a small core is one ingredient of the composition; every
verification claim must be read at its stated boundary, not at the
boundary a reader hopes for. \\
\textbf{Genode / Sculpt} & Sculpt describes a general-purpose OS built
from Genode's microkernel architecture with capability-based security,
sandboxed drivers, and VMs, at a 26.04 release and in day-to-day use by
its developers. The 26.04 release (2026-04-30) adds a platform-driver
gatekeeper that mediates which drivers may claim which devices,
IOMMU-based DMA confinement, and a ``black hole'' component that
swallows device-access requests from components without matching
capability grants (Labs 2026; Genode Foundation 2026). & seL4 proof
claims must not be transferred automatically to every Genode/Sculpt
configuration or to the applications it hosts. & The most active
capability-based desktop trajectory in this review --- the lineage of
(Levy 1984) --- worth watching as a structural direction rather than
treated as a deployable comparator for hostile desktop workloads
today. \\
\textbf{GrapheneOS} & Hardening atop Android's security model,
documented to keep even Google Play services inside the standard app
sandbox (GrapheneOS Foundation 2026a). The 2026 release documentation
adds per-application opt-outs from \texttt{hardened\_malloc} for
compatibility, a five-mode protection model for the USB-C pogo pins
against malicious accessories, and continued hardening of the
sandboxed-Play environment (GrapheneOS Foundation 2026b). & A mobile
comparator, not a desktop Linux replacement; its guarantees are scoped
to its platform. The per-app allocator opt-outs are themselves evidence
of compatibility pressure: even the most defaults-disciplined mobile
platform keeps an escape hatch. & The transferable lesson is shipping
useful applications with constrained authority by default --- mainstream
evidence that \breaktt{application\_confinement} with usable defaults is
an achievable product property, not an aspiration. \\
\textbf{Kali / Parrot} & Kali identifies penetration testing and
security auditing as its purpose (Linux 2026); Parrot documents
hardening measures while stating its core remains tuned for security and
forensics (Security 2026). & Offensive-tool availability is not evidence
of superior protection for the machine running those tools; both are
purpose-built workloads. & The label fallacy in its purest form: a
distribution's purpose statement constrains what its design should be
credited with, exactly as sec.~4 argues for
property-based rather than name-based assessment. \\
\bottomrule\noalign{}
\end{longtable}
\end{frame}

\begin{frame}[fragile,allowframebreaks]{What each comparator contributes
to the composition argument}
\protect\phantomsection\label{what-each-comparator-contributes-to-the-composition-argument}
The composition argument of this review holds that the security
properties that matter --- containment, authority limitation, integrity,
recoverability, update operations, supply-chain trust, usable delegation
--- are supplied by different components cooperating, not by any single
system name. Each comparator earns its place by isolating one term of
that argument.

\end{frame}

\begin{frame}[fragile,allowframebreaks]{What each comparator contributes
to the composition argument}
\textbf{Whonix contributes the networking boundary.} Its
gateway/workstation structure is the same design move as Qubes keeping
networking out of the privileged administrative domain
sec.~5, demonstrated in a privacy-focused deployment: the
component that meets the hostile network is a separate, disposable-trust
domain (W. Project 2026b). Its documented compromise-exposure warning is
equally instructive: anonymity infrastructure does not protect the
workstation's own credentials and browser data from a compromised
workstation (W. Project 2026c). The project's end-of-security-support
notice for Whonix 17 --- with deprecation notices arriving across early
and mid 2026 --- adds the maintenance dimension: a boundary design is
only as good as its supported line, and the project treats migration to
Whonix 18 as a user obligation, not an option (W. Project 2026a).
Anonymity and credential protection are distinct properties; confusing
\end{frame}

\begin{frame}[fragile,allowframebreaks]{What each comparator contributes
to the composition argument}
them produces designs that hide the operator's location while
surrendering the operator's secrets.

\textbf{Tails contributes state hygiene.} Amnesia is a deliberate
\breaktt{persistence\_recovery} strategy: the default system leaves
nothing to forensically clean and resets every session. The optional
Persistent Storage shows the honest tension in that design --- every
byte an operator chooses to persist is a byte that survives compromise
and must be protected by other means (T. Project 2026b). The 7.x series
on its Debian 13 base demonstrates that the hygiene can be sustained
across upgrades rather than sacrificed to them: automatic upgrades now
preserve the Persistent Storage, so the amnesic posture and an update
path no longer trade off (T. Project 2026c). The same release line
documents a Secure Boot certificate-expiry notice for users of the
earlier 7.7-era Secure Boot enablement --- a reminder that even an
\end{frame}

\begin{frame}[fragile,allowframebreaks]{What each comparator contributes
to the composition argument}
amnesic system has a boot chain with its own maintenance clock. For the
composition argument, Tails demonstrates that reducing what survives is
a legitimate alternative to defending what persists, and that the two
strategies must be chosen knowingly rather than mixed accidentally.

\textbf{OpenBSD contributes trusted-base discipline --- including
self-revision.} Privilege separation, W\^{}X, and secure-by-default
service configuration are documented properties of a codebase maintained
under an unusual security-first discipline (O. Project 2026d, 2026a).
The 2026 record sharpens the lesson. The errata process for releases 7.8
and 7.9 removed the pledge \texttt{tmppath} permission across a series
of errata and fixed a related namei race --- retracting a documented
mitigation whose race behavior made it a liability rather than an asset
(O. Project 2026b). Release 7.9 then introduced
\breaktt{\_\_pledge\_open(2)}, which narrows the window between a pledged
\end{frame}

\begin{frame}[fragile,allowframebreaks]{What each comparator contributes
to the composition argument}
path check and the actual open, and strengthened \texttt{pinsyscall} to
refuse \texttt{PT\_LOAD} mappings that would override pinned entry
points (O. Project 2026c). A trusted base that removes its own
mitigations through a public, dated errata trail demonstrates the
property that matters: the mitigation inventory is a maintained
engineering artifact, not an advertising list. The correct comparative
conclusion remains narrow: for narrowly scoped network services, OpenBSD
is a serious non-Linux comparator; for a browser-heavy development
workstation with GPU acceleration and AI tooling, the documented
evidence does not support a superiority claim. Its lessons --- small
trusted base, mitigations as defaults, mitigations subject to revision
--- transfer to the composition argument as design standards for
whichever platform actually runs the workload.

\end{frame}

\begin{frame}[fragile,allowframebreaks]{What each comparator contributes
to the composition argument}
\textbf{seL4 and Genode/Sculpt contribute the small-trusted-base and
capability-mediation trajectories.} seL4's formal verification
demonstrates that a practical microkernel can carry machine-checked
guarantees --- and its assumptions documentation shows how verification
claims must be scoped to boot, hardware, and DMA assumptions (seL4
Foundation 2026a; Klein et al. 2009). The 2026 record shows the frontier
still moving in two directions at once: the functional verification of
the MCS kernel completed on 64-bit RISC-V, and the first confidentiality
proof landed on AArch64 --- a non-interference property on the
architecture real phones and laptops run (seL4 Foundation 2026c);
Microkit 2.3.0 adds x86\_64 IOMMU support to the smaller deployment
path, acknowledging that even a small verified kernel must delegate
device containment to hardware (seL4 Foundation 2026b). Genode and
Sculpt demonstrate the same lineage's userland expression, and the 26.04
\end{frame}

\begin{frame}[fragile,allowframebreaks]{What each comparator contributes
to the composition argument}
release makes the capability argument concrete at the hardware boundary:
the platform-driver gatekeeper mediates which components may claim which
devices, the IOMMU confines what a driver's DMA can actually reach, and
the ``black hole'' component absorbs device-access requests from
components holding no matching capability --- a structural refusal of
the confused-deputy path rather than a policy warning about it (Labs
2026; Levy 1984). Capability mediation matters directly to agents: a
capability system is the classic answer to the confused-deputy problem
of an authority holder that cannot distinguish legitimate from
attacker-chosen requests (Hardy 1988). The caveat cuts both ways: proof
claims attach to configurations, not product names, and a capability
system hosting a compromised browser gains no automatic protection for
what the browser legitimately receives.

\end{frame}

\begin{frame}[fragile,allowframebreaks]{What each comparator contributes
to the composition argument}
\textbf{GrapheneOS contributes the shipped-defaults demonstration ---
and its limits.} A mainstream mobile platform that keeps even
first-party services inside the standard app sandbox shows that
constrained authority by default is a product decision, not a research
artifact (GrapheneOS Foundation 2026a). The 2026 documentation adds
three details that are as instructive as the headline claim:
per-application opt-outs from \texttt{hardened\_malloc} exist because
some applications fail under the hardened allocator, a five-mode
protection model now governs the USB-C pogo pins against malicious
accessories, and the sandboxed-Play environment continues to be hardened
in place (GrapheneOS Foundation 2026b). Each item is evidence of the
same design stance --- containment defaults with narrowly scoped,
documented escape hatches --- and of its cost: even the most
defaults-disciplined platform in this review maintains compatibility
\end{frame}

\begin{frame}[fragile,allowframebreaks]{What each comparator contributes
to the composition argument}
exceptions, which is exactly the maintenance burden the desktop
candidates of sec.~7 must budget for. The mobile
platform is not the answer for a workstation; it is evidence that the
property is shippable.

\textbf{Kali and Parrot contribute the negative control.} Their
documented purposes are offensive tooling and forensics (Linux 2026;
Security 2026), and neither project claims otherwise. Including them
makes the label fallacy vivid: installing an offensive toolkit does not
harden the machine that runs it, any more than a firewall vendor's
marketing hardens a desktop. Assessment must attach to documented
properties of the deployed system, not to the security connotations of a
name.

\end{frame}

\begin{frame}[fragile,allowframebreaks]{What each comparator contributes
to the composition argument}
Taken together, the comparators reinforce the review's central claim
rather than competing with it. None of them, alone, addresses both
failure paths of sec.~3: Whonix and Tails narrow
the exploitation exposure of networking and state but leave authorized
agent access untouched; seL4 and Genode narrow the trusted base without
mediating an agent's credentials; GrapheneOS constrains applications
without constraining what a persuaded agent may do with legitimate
access. The synthesis --- which components, composed, cover both paths
--- is the trust-domain architecture of
sec.~10; the scenario-level placement of these
comparators appears in sec.~16.
\end{frame}

\end{document}
